\section{Introduction}
\label{sec:intro}

Deep Reinforcement Learning (DRL) has achieved remarkable success, yet learning continuous control policies directly from visual inputs remains computationally demanding~\cite{mnih2015humanlevel}. Traditional actor-critic methods, such as Soft Actor-Critic (SAC)~\cite{haarnoja2018soft}, dominate these tasks due to their natural handling of continuous actions. However, value-based methods offer compelling off-policy sample efficiency. Recent advances in action-space discretization and network branching have made it possible to adapt Q-learning for continuous domains~\cite{tavakoli2019action}, providing a viable alternative to standard actor-critic setups.

To benchmark these capabilities, the DeepMind Control Suite~\cite{tassa2018deepmind} provides a rigorous physics-based simulation framework. In this study, we focus on learning movement strictly from visual pixels, utilizing environments that span from the moderately complex Hopper to the highly unstable 17-joint Humanoid. 

\begin{figure}[t]
    \centering
    \includegraphics[width=0.75\linewidth]{humanoid_pose.png}
    \caption{Example of the learned Humanoid walking posture using Rainbow, depicting the folded torso and stable gait.}
    \label{fig:humanoid_pose}
\end{figure}

This paper systematically compares a baseline Deep Q-Network (DQN) against the heavily augmented Rainbow DQN algorithm~\cite{hessel2018rainbow} on these high-dimensional tasks. We investigate how enhancements like distributional value modeling, Prioritized Experience Replay, and multi-step returns enable an agent to map raw pixel stacks to stable joint torques. By evaluating these algorithms across varying difficulty levels, we isolate the components strictly necessary to achieve complex bipedal balance using purely visual observations.

% --- TOP OF PAGE PLOT PLACEHOLDER ---
\begin{figure*}[t]
    \centering
    \includegraphics[width=\textwidth]{dqn_hopper_results.pdf}
    \caption{Training metrics of DQN for the Hopper environment (smoothed with an exponential moving average).}
    \label{fig:hopper_results}
\end{figure*}
% ------------------------------------
